\documentclass{situationist-preprint}

\title{Against Dialetheism}
\msid{LOG-03-a}
\msfamily{LOG-03}
\msdiscipline{Logic and Philosophy of Logic}
\mstype{Principal manuscript}
\mslean{lean/Philosophy/LOG/LOG\_03.lean}
\mskeywords{dialetheism; true contradictions; truth-value gaps; LP; strong Kleene; exclusion; liar paradox}

\begin{document}
\maketitle

\begin{abstract}
Dialetheism holds that some contradictions are true. Its case rests on the semantic paradoxes, Russell's set, instants of change and inconsistent law. We argue that there are no true contradictions and that each of these cases is better treated with a truth-value gap. The key observation, due in essence to Parsons, is that the glut theory and the gap theory assign exactly the same values to every sentence: the logic of paradox and strong Kleene logic share their truth tables and differ only in which value is designated. The dialetheist's cases therefore do not decide between gluts and gaps; what decides is a principle about designation. We defend one such principle, Exclusion: a sentence and its negation are not both true. We prove that Exclusion and the designation of truth leave only full truth designated, that no contradiction is then true, and that each of the dialetheist's cases takes the middle value and is neither true nor false. We add that the gap reading keeps modus ponens for the material conditional, which the glut reading loses. The results are formalized in Lean.
\end{abstract}

\section{The problem}

Dialetheism, the view that there are dialetheias, was given its modern logical form by \citet{Asenjo1966} and developed at length by \citet{Priest1979,Priest2006}, who argued that the semantic paradoxes are best taken at face value: the liar argument is sound, and the liar sentence is both true and false. The logic Priest proposed for this purpose, the logic of paradox LP, is paraconsistent, so that a true contradiction does not make every sentence true.

The case for dialetheism has two parts. The first is the argument from the paradoxes: the liar, Curry's sentence, Russell's set and their relatives seem to establish contradictions by reasoning whose every step is intuitively valid, and dialetheism lets one accept the reasoning \citep{Priest2006}. The second is a list of further cases in which, Priest argues, the world itself is inconsistent: the instant at which a person leaves a room, at which she is both in the room and not in it; legal systems whose statutes jointly require and forbid an act; and, for some authors, borderline cases of vague predicates \citep{Hyde1997}. \citet{Beall2009} defends a more restricted dialetheism, on which the only true contradictions are those generated by transparent truth, as spandrels of a device introduced for generalization.

The opposition is old and various. \citet{LittmannSimmons2004} collect the standard objections: that the dialetheist cannot express disagreement, that \enquote{just false} produces revenge, and that the dialetheist's logic is too weak for ordinary reasoning. \citet{Lewis1982} granted the usefulness of paraconsistent logics for reasoning about inconsistent information while denying that any contradiction is true. The dialetheist has replies to each.

We argue that there are no true contradictions, and that the dialetheist's cases are better treated with truth-value gaps. The argument does not dispute the data. It shows that the data are equally well fitted by gaps, and then argues that a principle of designation, which the gap theory satisfies and the glut theory violates, decides the matter. This agrees with LOG-01, which treats the liar as neither true nor false in a Łukasiewicz logic whose three-valued fragment is the one considered here.

\section{Framework}

\begin{definition}\label{d:vals}
Sentences built from atoms by negation, conjunction and disjunction take values in $\{0, 1, 2\}$, with
\[
\neg x = 2 - x, \qquad x \wedge y = \min(x, y), \qquad x \vee y = \max(x, y).
\]
These are the truth tables of strong Kleene logic K3 and of LP alike. The material conditional $A \supset B$ abbreviates $\neg A \vee B$.
\end{definition}

\begin{definition}\label{d:des}
A value is \emph{designated} if a sentence with that value is true. The gap reading designates only $2$, so that a sentence with value $1$ is neither true nor false; this is K3. The glut reading designates $1$ and $2$, so that a sentence with value $1$ is both true and false; this is LP.
\end{definition}

Definitions~\ref{d:vals} and~\ref{d:des} make precise the point of \citet{Parsons1990}: the gap theorist and the glut theorist agree on every valuation and disagree only about what the middle value means. We leave the set of designated values open and constrain it by two premises.

\begin{assumption}\label{a:norm}
\emph{Normality.} The value $2$ is designated.
\end{assumption}

\begin{assumption}\label{a:excl}
\emph{Exclusion.} For no value $x$ are both $x$ and $\neg x$ designated: a sentence and its negation are not both true.
\end{assumption}

\begin{assumption}\label{a:sym}
\emph{Symmetry.} In each of the dialetheist's cases there is a sentence $A$ such that $A$ and $\neg A$ have the same value.
\end{assumption}

Assumption~\ref{a:sym} is the dialetheist's own description of the cases. The liar is equivalent to its own negation; at the instant of leaving, the claim that the person is in the room and the claim that she is not have equal title; in the inconsistent statute, the act is required and forbidden by provisions of equal standing. We grant this description in full. Assumption~\ref{a:excl} is the premise the dialetheist denies, and Section~\ref{s:excl} defends it.

\section{The argument}

\begin{lemma}\label{l:des}
Given Assumptions~\ref{a:norm} and~\ref{a:excl}, a value is designated if and only if it is $2$.
\end{lemma}

\begin{proof}
The value $2$ is designated by Normality. If $0$ were designated, then $0$ and $\neg 0 = 2$ would both be designated, contrary to Exclusion. If $1$ were designated, then $1$ and $\neg 1 = 1$ would both be designated, again contrary to Exclusion.
\end{proof}

\begin{theorem}\label{t:nocontra}
Given Assumptions~\ref{a:norm} and~\ref{a:excl}, no contradiction is true: under every valuation, $A \wedge \neg A$ is undesignated.
\end{theorem}

\begin{proof}
The value of $A \wedge \neg A$ is $\min(x, 2-x)$, where $x$ is the value of $A$; this is at most $1$. By Lemma~\ref{l:des}, it is undesignated.
\end{proof}

\begin{theorem}\label{t:gaps}
Given Assumptions~\ref{a:norm}--\ref{a:sym}, in each of the dialetheist's cases the sentence $A$ has the value $1$, and neither $A$ nor $\neg A$ is true.
\end{theorem}

\begin{proof}
By Symmetry, $2 - x = x$, so $x = 1$. By Lemma~\ref{l:des}, neither $1$ nor $\neg 1 = 1$ is designated.
\end{proof}

Theorem~\ref{t:gaps} is the precise sense in which the dialetheist's cases are better treated with gaps: they keep the value the dialetheist assigns, and the principle of designation that rules out gluts converts that value into a gap. Two further facts bear on the comparison.

\begin{proposition}\label{p:mpgap}
On the gap reading, modus ponens for the material conditional preserves truth: if $A$ and $A \supset B$ are designated, so is $B$.
\end{proposition}

\begin{proof}
If $A$ has value $2$, then $\neg A$ has value $0$, and $\max(0, y) = 2$ only if $y = 2$.
\end{proof}

\begin{proposition}\label{p:mplp}
On the glut reading, modus ponens for the material conditional fails: if $A$ has value $1$ and $B$ value $0$, then $A$ and $A \supset B$ are designated and $B$ is not.
\end{proposition}

\begin{proof}
$A \supset B$ has value $\max(1, 0) = 1$.
\end{proof}

\begin{proposition}\label{p:dual}
$A$ entails $B$ on the gap reading if and only if $\neg B$ entails $\neg A$ on the glut reading.
\end{proposition}

\begin{proof}
On the glut reading, $\neg B$ is designated just when $B$ has value at most $1$, and likewise for $A$. So the right-hand side says that whenever $B$ falls short of $2$, $A$ falls short of $2$, which is the contrapositive of the left-hand side.
\end{proof}

Proposition~\ref{p:dual} is the duality between K3 and LP: neither theory is the stronger, and the choice between them must be made on grounds of what truth and negation are for.

\section{Why Exclusion}\label{s:excl}

The dialetheist rejects Assumption~\ref{a:excl}. We give three reasons for accepting it.

\emph{Rejection.} Asserting $\neg A$ is the standard way of rejecting $A$, and disagreement is expressed by one party asserting what the other denies. If $A$ and $\neg A$ can both be true, then asserting $\neg A$ does not rule $A$ out, and the dialetheist has no sentence whose assertion expresses rejection. \citet{Priest2006b} replies that rejection is a speech act, not the assertion of a negation, and that one can reject $A$ without asserting anything. But a speech act that rules $A$ out has a content, that $A$ is not true; if the content can be expressed it is a negation satisfying Exclusion, and if not, the dialetheist cannot say what she rejects.

\emph{Revenge.} The dialetheist must deny that \enquote{just false}, that is, false and not true, can be expressed so as to exclude truth; otherwise the sentence that says of itself that it is just false would be true only if it is not true, with no glut to absorb it. LOG-01 shows that the gap theorist faces the structurally parallel problem for \enquote{exactly true}, and answers it by proving that the corresponding operator is not definable. But the gap theorist loses an operator that ordinary reasoning does not use; the dialetheist loses the ordinary use of \enquote{not true}.

\emph{Detachment.} Proposition~\ref{p:mplp} shows that the glut reading loses modus ponens for the material conditional, and with it disjunctive syllogism. \citet{Priest2006} accepts this and adds a separate, relevant conditional that does detach. The gap reading keeps both inferences without supplementation (Proposition~\ref{p:mpgap}).

\section{The cases}

\emph{Semantic paradoxes.} The liar and its relatives take the middle value (Theorem~\ref{t:gaps}). LOG-01 extends the treatment to revenge sentences built from determinateness, and LOG-04 to Curry's sentence, which on our view proves nothing because contraction fails for the conditional. Curry's sentence threatens triviality, and dialetheism needs a contraction-free conditional to block it just as the gap theory does \citep{Beall2009}; gluts buy nothing here.

\emph{Russell's set.} The naive comprehension principle yields a set that is a member of itself if and only if it is not. The dialetheist accepts naive comprehension and the contradiction. On the gap treatment, membership in Russell's set is a sentence equivalent to its own negation and takes the middle value.

\emph{Instants of change.} At the instant a person leaves the room, Priest argues, she is both in and not in the room, since there is no reason to favor either. Symmetry is granted; it supports a gap at least as well as a glut. A classical treatment is also available: if the room is a closed region, she is in it at the last instant of contact.

\emph{Inconsistent law.} Courts resolve a conflict between statutes by a principle of precedence or declare the law indeterminate on the point; both responses treat the law as failing to settle the question, which is a gap, not a glut.

\emph{Vagueness.} \citet{Hyde1997} proposes that borderline cases are both $F$ and not $F$. The comparison of gluts with gaps is here beside the point, since we hold with MET-18 and LOG-15 that vague predicates are bivalent, with sharp but unknowable boundaries. If a non-classical treatment were wanted, Proposition~\ref{p:dual} shows that gaps would fit whatever data gluts fit.

\section{Objections and replies}

\begin{objection}
The argument assumes Exclusion, which is just the denial of dialetheism. It begs the question.
\end{objection}

\begin{reply}
Exclusion is the premise, and the dialetheist denies it; the argument is not meant to persuade someone who will give up nothing. What it shows is that the dialetheist's cases provide no evidence against Exclusion, since by Theorem~\ref{t:gaps} they are accommodated with Exclusion intact.
\end{reply}

\begin{objection}
The gap theorist cannot say that the liar is not true without saying something that, by her own lights, is not true. The dialetheist can say it truly.
\end{objection}

\begin{reply}
The gap theorist says that the liar is not determinately true, which is fully true (LOG-01, where the reply to Kripke's problem is given in detail). The dialetheist can say that the liar is not true, but she must add that it is also true, and her assertion therefore fails to communicate what the gap theorist's does.
\end{reply}

\begin{objection}
Paraconsistent logics are useful. Inconsistent databases, theories and belief sets are reasoned about without explosion. This supports dialetheism.
\end{objection}

\begin{reply}
It supports paraconsistency, not dialetheism. A paraconsistent consequence relation can be used to reason about inconsistent information without any of it being true, as \citet{Lewis1982} noted.
\end{reply}

\section{Formalization}

The module \nolinkurl{Philosophy/LOG/LOG_03.lean} formalizes Definitions~\ref{d:vals} and~\ref{d:des}, Lemma~\ref{l:des}, Theorems~\ref{t:nocontra} and~\ref{t:gaps}, and Propositions~\ref{p:mpgap}--\ref{p:dual}. The set of designated values is an axiom, \texttt{Designated}, constrained by \texttt{normality} and \texttt{exclusion} (Assumptions~\ref{a:norm} and~\ref{a:excl}); the cases and their values are axioms constrained by \texttt{symmetry} (Assumption~\ref{a:sym}). The main theorems are stated as follows.

\begin{lstlisting}
theorem no_true_contradiction (v : Nat → Nat) (A : Formula) :
    ¬ Designated (eval v (.conj A (.neg A)))
theorem cases_are_gaps (c : Case) :
    value c = 1 ∧ ¬ Designated (value c) ∧ ¬ Designated (2 - value c)
\end{lstlisting}

\begin{lstlisting}
'Philosophy.LOG_03.no_true_contradiction' depends on axioms:
  [propext, Classical.choice, Quot.sound,
   Philosophy.LOG_03.Designated, Philosophy.LOG_03.exclusion,
   Philosophy.LOG_03.normality]
'Philosophy.LOG_03.cases_are_gaps' depends on axioms:
  [propext, Classical.choice, Quot.sound, Philosophy.LOG_03.Case,
   Philosophy.LOG_03.Designated, Philosophy.LOG_03.exclusion,
   Philosophy.LOG_03.normality, Philosophy.LOG_03.symmetry,
   Philosophy.LOG_03.value, Philosophy.LOG_03.value_le]
\end{lstlisting}

Besides Lean's standard axioms and the primitives (\texttt{value\_le} bounds the values), the premises are \texttt{normality}, \texttt{exclusion} and \texttt{symmetry}. Theorem~\ref{t:nocontra} does not depend on Symmetry: that no contradiction is true follows from Normality and Exclusion alone, whatever the cases are.

\section{Conclusion}

The dialetheist correctly describes her cases as ones in which a sentence and its negation have equal standing. That description fits gaps exactly as well as gluts, since the two readings share every valuation. What decides between them is Exclusion, which preserves the use of negation to reject, avoids the revenge of \enquote{just false} and keeps detachment for the material conditional. Given Exclusion, no contradiction is true, and each of the dialetheist's cases is neither true nor false.

\begin{thebibliography}{99}\small
\bibitem[Asenjo(1966)]{Asenjo1966} Asenjo, F.~G. (1966). A calculus of antinomies. \emph{Notre Dame Journal of Formal Logic}, 7(1), 103--105.
\bibitem[Beall(2009)]{Beall2009} Beall, J.~C. (2009). \emph{Spandrels of Truth}. Oxford University Press.
\bibitem[Hyde(1997)]{Hyde1997} Hyde, D. (1997). From heaps and gaps to heaps of gluts. \emph{Mind}, 106(424), 641--660.
\bibitem[Lewis(1982)]{Lewis1982} Lewis, D. (1982). Logic for equivocators. \emph{Noûs}, 16(3), 431--441.
\bibitem[Littmann and Simmons(2004)]{LittmannSimmons2004} Littmann, G. and Simmons, K. (2004). A critique of dialetheism. In G. Priest, J.~C. Beall and B. Armour-Garb (eds.), \emph{The Law of Non-Contradiction}, Oxford University Press.
\bibitem[Parsons(1990)]{Parsons1990} Parsons, T. (1990). True contradictions. \emph{Canadian Journal of Philosophy}, 20(3), 335--353.
\bibitem[Priest(1979)]{Priest1979} Priest, G. (1979). The logic of paradox. \emph{Journal of Philosophical Logic}, 8(1), 219--241.
\bibitem[Priest(2006a)]{Priest2006} Priest, G. (2006a). \emph{In Contradiction: A Study of the Transconsistent} (2nd ed.). Oxford University Press.
\bibitem[Priest(2006b)]{Priest2006b} Priest, G. (2006b). \emph{Doubt Truth to Be a Liar}. Oxford University Press.
\end{thebibliography}

\end{document}
